\section{Introduction}
\label{sec:intro}
Immersive virtual reality (VR) is used in clinical rehabilitation, industrial and defense training and human-performance assessment, where the headset shapes perception, balance and spatial orientation rather than merely presenting information. Cybersickness---nausea, disorientation, oculomotor strain and, in severe cases, postural instability---is a long-standing side effect of the sensory conflict that VR creates~\cite{caserman2021cybersickness}. It can also be induced deliberately: an adversary who manipulates rendering parameters such as frame pacing, camera motion or field of view can make users sick or unsafe~\cite{tseng2022dark,odeleye2021framerate,kundu2025securing}. In both cases a detector that recognizes cybersickness during immersion, so that a mitigation can be triggered, is a safety function, and its two errors are not equally costly: a missed event leaves the user exposed to harm, whereas a false alarm causes a brief, recoverable interruption.

Detectors have progressed from post-exposure questionnaires to deep temporal and multimodal models of physiological, eye-tracking and gameplay data~\cite{dennison2016physio,islam2020automatic,islam2021fusion,liu2024cnneca,yalcin2024automatic,jeong2024precyse}. To encode unequal error costs, several authors have cast classification as a sequential decision problem and trained deep Q-networks with class-dependent rewards~\cite{lin2020imbalanced,yang2024drlimbalanced}. E-DDQN follows this line for cybersickness: a dueling Double DQN~\cite{vanhasselt2016double,wang2016dueling} whose reward penalizes a missed event five times more than a false alarm, trained with class-balanced replay and soft target updates.

Two questions must be answered before such a detector can be credited with a safety benefit. The first concerns evaluation. Gameplay telemetry is logged about once per second and self-reported discomfort changes rarely within a session, so any split that places seconds of the same session or participant in both training and test data rewards recognition of the session rather than detection of cybersickness. Oversampling before splitting is a known source of such leakage~\cite{kapoor2023leakage}, and published results on the gameplay dataset used here were obtained by splitting samples rather than participants~\cite{kundu2022truvr}. An earlier, unpublished version of the present study reported \pct{96.88} accuracy for E-DDQN under an oversampling-before-splitting protocol. The second question concerns the formulation. It is rarely made explicit what an asymmetric reward adds to classical cost-sensitive learning~\cite{elkan2001cost}, and RL detectors are seldom compared with cost-weighted supervised models tuned and thresholded under the same rules.

We therefore address four research questions on the public UFFCSData dataset~\cite{porcino2019uffcsdata}:
\begin{itemize}
\item[RQ1] How much does the split protocol change the measured performance?
\item[RQ2] What does the asymmetric reward add, formally, to cost-sensitive classification?
\item[RQ3] Under leakage-free evaluation, does E-DDQN outperform a standard DQN, cost-weighted supervised models and a minimal exposure-time baseline, and does any model meet a safety-prioritized operating target?
\item[RQ4] How sensitive are the detectors to bounded manipulation of the telemetry channels that an adversary could influence?
\end{itemize}

The contributions are:
\begin{enumerate}
\item \emph{Leakage audit.} We show that oversampling before a random row-level split places exact copies of \LeakDupPct\,\% of the positive test rows in the training data and lets an untuned classifier reach ROC-AUC \LeakOldRocAuc, whereas a session-grouped split of the same features gives \LeakGrpRocAuc{} (Section~\ref{sec:leak}).
\item \emph{Formal link between RL and cost-sensitive learning.} We prove that, on logged data, the alarm cannot change future observations, so detection is a contextual bandit, and that the 5:1 reward induces exactly the Bayes decision rule $P(y=1\mid s)\ge1/6$ of 5:1 cost-sensitive classification (Section~\ref{sec:formulation}).
\item \emph{A matched, leakage-free benchmark.} We compare E-DDQN with a standard DQN, a contextual-bandit variant and four cost-weighted supervised models under a nested, session-grouped protocol fixed before testing, with identical tuning budgets, checkpoint and threshold rules and paired statistics; two exposure-time baselines were added post hoc under the same protocol (Sections~\ref{sec:groups}--\ref{sec:protocol}).
\item \emph{Findings.} E-DDQN does not outperform the standard DQN and has a higher cost than a cost-weighted random forest; elapsed time alone ranks seconds as well as any model; and no model combines high recall with a controlled false-alarm rate (Section~\ref{sec:results}).
\item \emph{Diagnosis and stress tests.} Bootstrapping and validation-based selection explain E-DDQN's under-performance (Section~\ref{sec:why}); bounded telemetry perturbations show that all detectors rely mainly on exposure time and questionnaire answers (Section~\ref{sec:attack}).
\item \emph{Reproducibility.} Code, the fixed protocol and per-sample predictions are released so that every number in the paper can be regenerated from the raw data.
\end{enumerate}

Section~\ref{sec:related} reviews related work, Section~\ref{sec:method} presents the methods, Section~\ref{sec:setup} the setup, Section~\ref{sec:results} the results and Section~\ref{sec:discussion} the discussion; Section~\ref{sec:conclusion} concludes.
